% BFZI P0, 2026-11-04. Source-only derivation and validation contract.
\section{P0：锁定磁准静态物理域、连续闭合源项与有限直导线参考}
\label{sec:p0-physical-source}

\subsection{单一物理场景与相量符号}
在本节首先只讨论静止、线性、$\mu=\mu_0$、
海水标量电导率 $\sigma_s>0$、
两只互不相交的理想绝缘实心封装及
BFZI 圆形--圆角方形\textbf{单根连续螺旋}。
每个端口的电流沿螺旋外端到内端，
再沿显式规定的封装内回流折线回到外端，
形成相同方向约定下的闭合曲线 $\Gamma_p$。
正频相量为 $e^{+\jj\omega t}$，端口电流使用峰值。
圆形情形是角部连续绕行与直边长度为零的
合法极限，不重新引入旧超椭圆族。

海水外部真实区域与人工有限计算区域不同：
\[
\Omega_\infty=\RR^3\setminus(\overline B_1\cup\overline B_2),
\qquad
\Omega_R=B_R\setminus(\overline B_1\cup\overline B_2).
\]
封装界面满足 $\bm n\cdot\bm J_s=0$，
而计算球的相同零法向条件只是
\emph{人为的有限域电流条件}。
在未对人工球和无限海水外区实现非局部相容耦合前，
不得称 $\Omega_R$ 的离散解为连续无限域真值。

\subsection{涡流近似的必要尺度核查}
海水中忽略位移电流并不等于自动满足
Maxwell--MQS 的所有适用条件：
至少需检验
\[
\eta_{\mathrm D}
=\frac{\omega\epsilon_s}{\sigma_s}\ll1,
\qquad
\eta_{\mathrm wave}
=k_0D=\frac{\omega D}{c_0}\ll1,
\qquad
\delta_s=\sqrt{\frac{2}{\omega\mu_0\sigma_s}} .
\]
还要检查真实介质波数与装置尺度、
源电流的连续性、界面极化/介电损耗是否可忽略。
Ammari--Buffa--N\'ed\'elec~\cite{ammaribuffanedelec2000}
强调涡流近似的高阶准确性依赖
外加源所满足的额外兼容条件。
例如 $f=200\,\mathrm{kHz}$、
$\sigma_s=5\,\mathrm{S/m}$、
$D=0.52\,\mathrm m$ 时
$\delta_s\approx0.503\,\mathrm m$，
$k_0D\approx2.18\times10^{-3}$；
\emph{仅作例示}若设海水相对介电常数为 $80$，
则 $\eta_{\mathrm D}\approx1.78\times10^{-4}$。
当前海水模型尚无实际的色散介电输入，
这些示例数字\textbf{不提供任意海水工况的电磁误差界}。

\subsection{同一闭合路径的理想线源核}
将外部指定单位电流源记作
\[
\bm C_p(\bm x)=\frac1{4\pi}
\oint_{\Gamma_p}\frac{\mathrm d\bm\ell'}{|\bm x-\bm r'|},
\qquad
\bm A_{\mathrm{src},p}=\mu_0\bm C_p,\qquad
\bm B_{\mathrm{src},p}=\mu_0\Curl\bm C_p .
\]
它只在实际海水域中评价，避免零半径导线上的
奇异核。闭合曲线积分与
\emph{有限节点点线元加权和}不是同一个分布：
\[
\bm C_{p,h}(\bm x)=\frac1{4\pi}
\sum_k\frac{\bm h_k}{|\bm x-\bm r_k|}.
\]
两者在光滑、离线距离有正下界的查询区域可通过
积分阶次收敛相互逼近。
虽然 $\bm B_{p,h}=\mu_0\Curl\bm C_{p,h}$
对上述同一核严格成立，
且 $\sum_k\bm h_k\approx0$，
仍不自动使有限点线元分布成为
严格无散度的物理闭合回路电流。
必须独立检查盒壁附近的
$\bm C_{p,h}$、$\bm B_{p,h}$ 求积误差。

\subsection{不依赖源 Gauss 节点的精确有限直导线核}
对起止点 $\bm a\to\bm b$，
令 $L=|\bm b-\bm a|$、
$\bm t=(\bm b-\bm a)/L$、
$\bm R_1=\bm x-\bm a$、
$\bm R_2=\bm x-\bm b$、
$r_1=|\bm R_1|$、$r_2=|\bm R_2|$。
从线积分直接计算得到
\[
\boxed{
\bm C_{\mathrm{seg}}(\bm x)
=\frac{\bm t}{4\pi}
\log\frac{r_1+r_2+L}{r_1+r_2-L}
}
\]
和
\[
\boxed{
\bm B_{\mathrm{seg}}(\bm x)
=\frac{\mu_0}{4\pi}
\frac{(\bm R_1\times\bm R_2)(r_1+r_2)}
{r_1r_2(r_1r_2+\bm R_1\cdot\bm R_2)} .
}
\]
仅在不位于理想导线的场点使用。
这两个公式无近奇异 Gauss 点问题，
可以作为 BFZI 每根\textbf{真实直边}
及每段明确回流线的独立源场参考。
BFZI 角部仍然是连续曲线：
固定角部几何和真实有向线元，
仅对角部使用高阶 Gauss 收敛，不得
用圆弧中心拟合、独立闭合匝或假定匝同心
来替代实际 BFZI 公式。

代码 \texttt{sdfmpneo\_bfzi/source\_reference.py}
提供 \texttt{straight\_segment\_field} 与
\texttt{continuous\_port\_field}。
首轮 P0 本机 \texttt{a50ed62} 的
75 项测试通过；独立角段 24/32 阶积分一致至机器精度，
但旧生产源的 4 阶整路径 Gauss
$\bm C/\bm B$ 探针误差分别约
$0.253\%$ 和 $1.174\%$，
12 阶时分别约 $0.000118\%$ 和 $0.00380\%$。
因此旧 4 阶源场\textbf{未达到}预先规定的
$10^{-3}$ 门槛。

随后已经据此修正\texttt{LineSource}：
实际端口源场 $\bm C$ 与 $\bm B$ 的直边
和全部回流段使用上述解析有限直导线核，
曲角才采用该阶次的 Gauss 节点。
完整的旧线元节点仍独立保留用于导线长度、
铜线热源和空气 Neumann 互感近似，
避免错误地把上述源场修正称为已修复铜线
邻近效应或完整空气自感。
源场两种读出继续使用同一
$\bm C$ 的解析旋度，即
$\bm B=\mu_0\Curl\bm C$；
本机在代码 \texttt{ce086b1} 上最新
77 项测试通过，并在两端口各 28 个海水探针上测得：
新 4 阶 $\bm C/\bm B$ 的相对误差分别为
$2.24\times10^{-5}$ 与 $1.95\times10^{-4}$，
原全 Gauss 4 阶则为
$2.53\times10^{-3}$ 与 $1.17\times10^{-2}$。
故\textbf{固定示例上的源场已通过预定 $10^{-3}$ 门槛}，
不能将该点集结果外推全部几何参数域或
真实有限截面导线近邻电磁分布。

\paragraph{闭合回路的 Coulomb 规范连续性检查}
完整的光滑闭合源曲线在导线外满足
\[
\Div\bm C_p(\bm x)
=-\frac1{4\pi}\oint_{\Gamma_p}
\frac{\mathrm d\bm\ell'\cdot(\bm x-\bm r')}
{|\bm x-\bm r'|^3}
=-\frac1{4\pi}\oint_{\Gamma_p}
\mathrm d_{\ell'}\frac1{|\bm x-\bm r'|}
=0.
\]
对任何直线段 $\bm a\to\bm b$，
这项的精确贡献为
\[
\Div\bm C_{\rm seg}(\bm x)=
\frac1{4\pi}\left(
\frac1{|\bm x-\bm a|}
-\frac1{|\bm x-\bm b|}\right).
\]
沿物理闭合路径各段端点项望远镜式抵消；
仅曲角用有限 Gauss 时残留可测数值偏差。
现在实现的
\texttt{LineSource.divergence\_kernel}
直接积分这些解析端点项和曲角点线元的导数核，
而不是通过 $B=\mu_0\Curl C$ 自我验证。
同时新增盒面、棱、角外点的局部源场误差。
本项扩展诊断尚待本机运行，不作为新的
全 Maxwell 认证或非磁性有限截面适用域证明。


\subsection{P0 连续 BFZI 几何的角部正向性约束}
支撑函数角部的有向切线在两种绕向下都满足
\[
\frac{\mathrm d\bm r_{m,k}}{\mathrm du}
=\sigma_w\!\left[\varrho(\eta)-
2\sigma_w(1-\alpha)\beta
(\bm e_m\cdot\bm t(\psi_m))\right]\bm t(\psi_m),
\qquad
\sigma_w(\bm e_m\cdot\bm t)=\cos u-\sin u\in[-1,1].
\]
所以，\emph{只检查 $\varrho(\eta)>0$ 不足以防止角部
切线消失或反向}。
目前的 BFZI 构造器加上一个直接有理论依据的
保守充分条件（并非必要条件）：
\[
\boxed{\;
\varrho_{\min}=\varrho_0-\alpha sN
>2(1-\alpha)\beta,\qquad\beta=s/(2\pi).\;}
\]
若违反该条件，
不得把仍为正数的角半径误当成一根正向无折返的
圆角方形连续导线。这避免错误源几何进入求积和
后续海水教师。有限半径导线的自接触/弯曲半径
尚需更强的几何条件，不能从角切线正向性单独推出。

\subsection{严格闭合折线参考与连续曲角近似误差}
本机 \texttt{325d47b} 的 P0 日志（79 项测试通过）
将双端口的查询从 28 点扩至各 84 点，
增加真实盒面/棱/角附近的海水点：
4 阶源的 $\bm C,\bm B$ 总体误差分别
$7.28\times10^{-5},4.27\times10^{-4}$，
边角 $\bm B$ 误差 $7.48\times10^{-4}$；
而该源的归一化 $\Div\bm C$ 仍为
$1.06\times10^{-4}$，从 6、8、12 阶
分别下降至
$3.31\times10^{-6},1.02\times10^{-7},8.86\times10^{-11}$。
这个结果支持曲角 Gauss 残差逐级收敛，
\emph{但它不是严格满足无散度源条件的离散源}。

因此另行构造真正闭合的折线
$\bm v_0,\ldots,\bm v_M=\bm v_0$，
每段为从解析 BFZI 曲角上选取的相邻真实位置
形成的弦、完整解析直边或原回流折线。
所有段仍使用未经软化的有限直线精确核；
任意离散阶数、任意导线外查询位置，
其散度严格望远镜式抵消：
\[
\Div\bm C_{\rm poly}(\bm x)
=\frac1{4\pi}\sum_{j=0}^{M-1}
\left(\frac1{|\bm x-\bm v_j|}
-\frac1{|\bm x-\bm v_{j+1}|}\right)
=0.
\]
该性质\textbf{完全独立于曲角 Gauss 求积}；
且曲角只在正切线的原 BFZI 路径上细分，
所以不引入任何历史超椭圆或独立闭合同心圆。
然而折线与实际连续曲线并不相同，
要使 $\bm C_{\rm poly},\bm B_{\rm poly}$
逼近连续源仍必须做角部弦段
16/32/64 个面板的\emph{几何收敛}比较。
闭合性是物理合法性，场精度是另一个量。
代码接口为
\texttt{polygonal\_port\_vertices} 和
\texttt{polygonal\_port\_field}。
其新增实验与测试尚待本机验证；
暂不切换默认生产源算法。

\subsection{P0 验收和与 P1 的严格分界}
P0 必须在固定两端口位姿、
真实连续中心线和完整回流下执行：
\begin{enumerate}
\item 直边、角部及相邻匝真实首尾连接；
线圈外端到内端再沿回流回到外端；
旋转变换和源向量方向符合右手坐标。
\item 在与真实直线源远离的场点，解析
$\bm C_{\mathrm{seg}},\bm B_{\mathrm{seg}}$
与足够高阶的独立 Gauss 源积分相等；
检查刚体协变和
$\bm B_{\mathrm{src}}=\mu_0\Curl\bm C_{\mathrm{src}}$。
\item 圆形与圆角方形均以独立参考阶次
\texttt{corner\_order=24/32} 比较场；
固定盒壁、边缘、角点和远场查询点，
报告生产 \texttt{order=3/4/6/8/12}
的源场相对误差，不以端口能量闭合替代。
\item 对每个目标介质/频率重新核查
$\omega\epsilon_s/\sigma_s$、$k_0D$、
皮肤深度和内部导体材料近似。
本节不认证超出线性均匀海水的情形。
\item P0 的正确闭合源只是
下一阶段原 Coulomb 磁电感积分的
\emph{已冻结外部激励}；在磁能
$L_{mn}$ 尚未独立收敛前，
不能宣称海水 $J$、$q$、阻抗或神经教师达到真值。
\end{enumerate}
本机入口是
\texttt{python run.py --mode bfzi-p0-source --with-tests}，
将只输出源项与尺度的
\texttt{p0\_physics/p0\_reference/p0\_source}，
不启用 GitHub Actions。
原解析源修正已在本机 77 项测试中通过，新增闭合\texttt{div C}
与盒棱/角源场扩展诊断尚待本机复测。
当前\textbf{P0 固定示例源精度通过，但参数域尚未全部验收，G1=FAIL}。
